\section{A broad side with a small polynomial density cap}
\label{sec:very-small-polynomial-width-on-the-broad-side-polynomial-defect}

\begin{lemma}[A broad side with a small polynomial density cap]
\label{lem:small-polynomial-broad-side}
Work in the counterexample sequence fixed above. There is a universal \(D_*>0\) with the following property. Let \(0<\gamma<1,\ p_0>0\) be fixed. Suppose a finite random tag of law \(\Lambda\), with pairs \((\mu_i,\nu_i)\) on \(X,Y\), satisfies
\[
 N\mathbb E_\Lambda\mu_i,\ N\mathbb E_\Lambda\nu_i\le K,\qquad
 \operatorname{width}(\mu_i)\le n^\gamma,\qquad N\max\nu_i\le n^{D_*},
\]
where \(K\) is a constant and both average inequalities hold pointwise. In one color \(G\) assume on every support \(y\in\operatorname{supp}\nu_i\), \(d_G(\mu_i,y)\ge1-\epsilon,\ \epsilon=n^{-p_0}\). Then one of the colors contains the cube for large \(n\).
\end{lemma}

\begin{proof}
The preliminary reduction either produces a cube in the opposite color
or supplies pairs of parent labels that both see a common first-side law.
In the latter case we construct first-side tuples hitting those parents,
then use their observations to assign the odd roles. The posterior of a
whole tuple after its odd neighbors are assigned will supply the even
row for Hall's theorem. This follows the parent and height construction
of Lemma~\ref{lem:constant-broad-high-neighborhood}, with two changes:
each tuple carries a tag that must be deleted along with its entries,
and a high odd row redraws an actual parent rather than a hidden vector.
The constant \(D_*\) is universal; only subsequent auxiliary choices
and the sufficiently large dimension threshold may depend on the fixed
input parameters.

\subsection{Local base histories and hidden tags}
\label{sec:local-base-histories-and-hidden-tags}

The second-side mixture and its tag posterior are
\[
 \Pi(y)=\sum_i\Lambda(i)\nu_i(y),\qquad
 \eta_y(i)=\frac{\Lambda(i)\nu_i(y)}{\Pi(y)},\qquad
 B_y=\sum_i\eta_y(i)\mu_i
\]
when \(\Pi(y)>0\). Thus \(B_y\) is the first-side law obtained by
observing \(y\) in the second-side mixture. Restrict \(\Pi\) and normalize
to obtain \(\Pi'\) on \(\{y:N\Pi(y)\ge n^{-D_*}\}\). The discarded mass
is at most \(n^{-D_*}\), and \(\Pi'\le O(K)/N\).

Suppose first that two independent \(\Pi'\)-draws satisfy
\(d_G(B_y,y')<c_0=.01\) with probability at least \(.1\).
A set of first draws of positive \(\Pi'\)-mass then has a positive
\(\Pi'\)-mass of such partners. Each partner set contains \(\Omega(N)\)
labels because \(\Pi'\le O(K)/N\). The laws \(B_y\) have width at most
\(n^\gamma\), and their average over these first draws is at most
\(O(1)/N\): indeed \(\mathbb E_{y\sim\Pi}\eta_y=\Lambda\).
Their degrees to the partners in the other color exceed \(.99\).
Lemma~\ref{lem:constant-broad-high-neighborhood} therefore gives a cube
in that color.

In the remaining case the directed bad-pair probability is less than
\(.1\). A union bound for the two directions shows that the relation
\[
 y\asymp y'\quad\Longleftrightarrow\quad
 d_G(B_y,y'),d_G(B_{y'},y)\ge c_0
\]
has \((\Pi'\otimes\Pi')\)-mass greater than \(.8\). In particular a
positive \(\Pi'\)-mass of labels has partner mass at least \(.1\).
Draw \(V_0\) from \(\Pi'\) restricted to those labels. Given \(V_0\),
draw the candidates \(A_w\) independently from
\(\Pi'|_{\{y:y\asymp V_0\}}\), with bins \(w\) defined below.
Both the initial prior and every conditional candidate prior are bounded
by constant multiples of \(\Pi\).

Use the coarse and fine chunks of
Lemma~\ref{lem:constant-broad-high-neighborhood}: there are \(d_c=300\)
coarse chunks of length \(\lfloor n^{1/5}\rfloor\), and
\(m=\lceil n^\alpha\rceil\) fine chunks of odd length asymptotic to
\(n^{.3}\). Their majority signs, exactly flippable coordinates, and
number of counts within \(R_f=5.5\) of mid-weight are denoted by
\(t,F,j\). Here use \(J=\lfloor m^{.04}\rfloor\), and call \(j\le J\)
low. The signs remain uniform on either cube parity, while
\(\Pr(j\ge q)\le n^{-.13q}\) and the coarse boundary fraction is at most
\(n^{-.05}\).

The coarse key of a site \(x\) is \(h_x=(w,\mathrm{int/bdy})\).
Define its named primary variable by
\[
 P_{(w,\mathrm{int})}=A_w,\qquad
 P_{(w,\mathrm{bdy})}=V_0,
\]
and let \(P_h^*\) be the other member of \(\{A_w,V_0\}\).
Two keys have matching primaries when these are the same random variable,
not merely when their realized labels agree. Give the coarse keys a
symmetric relation connecting each key to itself, the two keys at a bin,
and boundary keys at adjacent bins; let \(C(h)\) include all these
neighbors and \(h\) itself. A coarse flip stays in this relation: if it
changes the bin, both endpoints are boundary keys. Within \(C(h)\), only
the opposite key at the same bin can have a nonmatching primary.

Given the parents, draw independent base tags \(I_h\) from
\(\eta_{P_h}\) restricted to
\(d_G(\mu_i,P_h^*)\ge c_1=c_0/2\).
The mean of this degree under \(\eta_{P_h}\) is at least \(c_0\), so
the restriction has mass bounded below by a positive constant.
On the parent support,
\[
 \eta_y(i)=\Lambda(i)\frac{\nu_i(y)}{\Pi(y)}
 \le n^{2D_*}\Lambda(i).
\]
Thus the conditional base-tag law is at most
\(O(n^{2D_*})\Lambda\).

Next, given the base \((V_0,A,I)\), draw independent hidden scalars
\(Z_\ell\), indexed by \(\ell=(h,t)\). Their laws
\(\pi_\ell=\pi_h\) are parent posteriors in this base model:

\begin{itemize}
\item at \(h=(w,\mathrm{int})\), redraw \(A_w\) given \(V_0\) and the tags \(I_s,\ s\in C(h)\);
\item at \(h=(w,\mathrm{bdy})\), redraw \(V_0\) given the candidates \(A_u\) at bins of \(C(h)\), and the tags \(I_s,\ s\in C(h)\).
\end{itemize}
Only low odd positions will use a hidden scalar \(Z_\ell\) as their target
label. A high position will instead use its actual named primary. Thus
there are no high hidden vectors or stream pools in this construction.

An even position \(x\) of coarse key \(h=h_x\) uses a type \(\beta\) carrying \(h\), low/high and the following observation list \(S\) (also retain \(j\) for low types):
\begin{itemize}
\item low: all \((s,t),\ s\in C(h)\), and \((h,t^a),\ a\in F\) with \(t^a\) the one-sign flip;
\item high: just \((h,t)\) when \(j=J+1\), and an empty list otherwise.
\end{itemize}
Write \(u_\beta=j+1\) for low, 1 for high. All adjacent low odd keys are covered (the high-to-low transition can only be at the paired fringe with sign unchanged).

For every potential center ID we will generate one tuple per type \(\beta\), independent across ID/type pairs given history through \(Z\), consisting of one new tag \(i\) and \(k=\lceil\kappa J\log n\rceil\) labels on \(X\). The tag has a restricted posterior law \(T_\beta\) defined below. Given this tag, draw iid from \(\mu_i\) conditioned individually to hit:
\begin{itemize}
\item \(P_h\) and all \(Z_S\) for low;
\item \(P_h,P_h^*\) and all \(Z_S\) for high.
\end{itemize}
These requirements include the actual primaries needed by adjacent high
odd roles. At a low-to-high transition the coarse key is unchanged.
Step~2 will prove that the common-hit mass is at least \(c_1/2\), so the
sampler is defined on its success gates.

The raw generation order is
\[
 V_0\ \longrightarrow\ (A_w,I_{(w,\mathrm{int})},I_{(w,\mathrm{bdy})})_w
 \ \longrightarrow\ (Z_\ell)_\ell
 \ \longrightarrow\ \text{center tuples}.
\]
Given \(V_0\), the bin variables in the second stage are independent;
given the base, so are the hidden scalars; given both, so are the tuples
at distinct ID/type pairs. Generate tuples also for IDs whose centers are
not present, independently of center positions, activations and ties.
On undefined branches use fixed local fallbacks respecting the bin, sign
and ID symmetries. Whenever a defining sampling formula is valid, use it
without checking numerical tests elsewhere. Later avoidance changes the
stage law, but does not redefine any raw posterior formula.

A tuple observation is the entire primitive pair consisting of the new
center tag and its \(k\) label entries. Deleting a tuple will remove that
whole observation. When the pair is varied in the even reconstruction,
all choices and observations derived from it are recomputed. The old base
tag \(I_h\) remains in the history: the fresh tag sampled from
\(T_\beta\) does not replace it. Likewise equal realized values of two
primary or hidden variables do not identify their roles in a conditioning
formula.

Parameters can be chosen as follows. Take the small fixed \(\kappa>0\) so small that \(\kappa(1+\log(4/c_1))<.001\). Use very small successive constants \(d_2,\delta_2,d_1,\delta_1,d_0,D_*>0\), each sufficiently small in terms of previous ones and the absolute coarse dimension (in particular \(d_2\ll\kappa\)). Choose \(2D_*<d_0\), so the fixed multiplier in the base-tag bound is
absorbed and that bound is \(n^{d_0}\Lambda\) for large \(n\).
If \(C\) is an absolute constant with \(|S|\le C u_\beta\), require
\(C d_1+\delta_2+d_0<d_2\), with slack; also take
\(d_0,\delta_1\ll d_1\). Finally \(\alpha>0\) is sufficiently small in terms of all these constants and \(p_0\). The target fan is \(T=\lceil m^{.001}\rceil\). Constants hidden in pattern counts below are absolute for the fixed coarse dimension; in particular choosing \(D_*\) does not require \(K,\gamma,p_0\). We detail the needed predictive tests.

\paragraph{Step 1: deleting one base tag from a hidden-key posterior.}
For \(\ell=(h,t)\), let \(\pi_{\ell,-s}\) be the defining parent
posterior for \(\pi_\ell\) with the observation \(I_s\) omitted and all
its other listed observations retained. We require
\[
 N\max\pi_\ell\le n^{d_1},\qquad
 \pi_\ell\le n^{d_1}\pi_{\ell,-s}\quad(s\in C(h)).
\]
To estimate a deletion failure, fix the retained observations of this
parent posterior. Given any supported candidate parent value, the
incoming tag has likelihood at most \(n^{d_0}\Lambda(i)\), by
conditional independence of the base tags and their uniform domination.
Its predictive density relative to \(\Lambda\) is smaller than
\(n^{-\delta_1}\) with probability at most \(n^{-\delta_1}\).
Outside that event Bayes' formula bounds the posterior multiplier by
\(n^{d_0+\delta_1}\le n^{d_1}\), simultaneously at every parent label.

For the absolute cap, use all observations in the parent posterior.
In the boundary case reference the bounded list of candidate labels to
uniform laws; reference all incoming tags to \(\Lambda\). The parent
and conditional candidate priors have normalized caps \(O(1)\), and
only boundedly many tag factors occur. The same predictive-denominator
calculation, with \(d_0,\delta_1\ll d_1\), proves the absolute test.
Each Step~1 failure has probability \(O(n^{-\delta_1})\) in its raw
parent experiment.

\paragraph{Step 2: reconstructing the tag of a center tuple.} For a type \(\beta\) at \(h\), fix the base except \(I_h\).
Let \(T_0\) be the original conditional law used to generate that base
tag from its two parents. We first replace \(I_h\) by a hypothetical
value \(i\) to define a posterior law on tags. Given the resulting
\(Z_S\), a center later samples a fresh tag from this law. The posterior
formula itself does not consult the actual value of \(I_h\) again. Restrict alternative tags \(i\) to \(G_\beta\) requiring the step-1 pointwise ratio comparisons for deletion of \(I_h\) at the keys in \(S\), computed with its value \(i\) (laws compared on all labels, so \(G_\beta\) does not read realized \(Z\)'s). Write
\[
 m_S=\int 1_{G_\beta}(i)\prod_{\ell\in S}
 \frac{\pi_\ell(Z_\ell;\,i)}{\pi_{\ell,-h}(Z_\ell)}\,dT_0(i).
\]
Notation \(m_{S-\ell}\) drops one factor with the same gate. Use their normalized integrands (against \(T_0\)) as \(T_\beta,T_{\beta,-\ell}\). Enforce positivity where needed and \(m_S,\ m_S/m_{S-\ell}\ge n^{-\delta_2 u_\beta}\). Conditional on the base with \(I_h\) omitted, the reference law of
\(Z_S\) is \(\prod_{\ell\in S}\pi_{\ell,-h}\).
The true-event data subdensity relative to this reference is \(m_S\).
Both \(m_S\) and \(m_{S-\ell}\) integrate to at most one, so for
\(a=n^{-\delta_2u_\beta}\),
\[
 \int_{m_S<a}m_S\le a,
 \qquad
 \int_{m_S<a m_{S-\ell}}m_S\le a\int m_{S-\ell}\le a.
\]
These are bounds on the joint event with true \(I_h\in G_\beta\).
Reference atoms of mass zero may be omitted on that gate. Division of the
normalized integrands now gives
\[
 T_\beta\le n^{d_2 u_\beta}\Lambda,\qquad T_\beta\le n^{d_2 u_\beta}T_{\beta,-\ell}.
\]
Indeed the absolute ratio to \(\Lambda\) is at most
\(n^{d_0+d_1|S|+\delta_2u_\beta}\), and the ratio to the deleted
posterior is at most \(n^{d_1+\delta_2u_\beta}\). The stated order of
constants places both below \(n^{d_2u_\beta}\). The same bounds hold at alternative observations satisfying the gates. Also every tag in \(T_\beta\)'s support sees \(P_h\) with degree \(\ge1-\epsilon\), \(P_h^*\) with degree \(\ge c_1\), and each \(Z_S\) with degree \(\ge1-\epsilon\) if its primary matches, or \(\ge c_1\) if not (degree here into \(\mu_i\)). Indeed positive likelihood for \(Z_{(h',t')}\) entails a positive case in the base model of \(P_{h'}\) having that value with \(I_h=i\), since \(h\in C(h')\). In that case the base tag at \(h\) was drawn to see that value with degree \(\ge1-\epsilon\) if its primary is the same named variable \(P_h=P_{h'}\), or degree at least \(c_1\) in the crossing case \(P_{h'}=P_h^*\). For a low tuple, the only required column that can have degree merely
\(c_1\) is the single \(Z\)-key with nonmatching primary. For a high
tuple it is \(P_h^*\); its optional \(Z\)-key has matching primary.
Every other required column has defect at most \(\epsilon\).
Consequently their common neighborhood has \(\mu_i\)-mass at least
\(c_1-O(m\epsilon)\ge c_1/2\). This proves that the tuple sampler is
defined on all the stated gates.

\subsection{Odd observation configurations and deletions}
\label{sec:odd-observation-configurations-and-deletions}

A state retains the residual bits and the fine count fields and total
\(j\) of the preceding lemma, with counts at distances \(R_f,R_f+1\)
merged on each side of mid-weight. On the coarse chunks it retains only
the bin vector, encoded by one-hot fields, and the interior/boundary flag.
The ambient encoded cube still has dimension \((1+o(1))n\).
Unlike the earlier state, this one need not determine parity, because
it omits exact coarse counts.

\Needspace{7\baselineskip}
Let \(q_{\rm st}\) be this quotient map, and set
\(\mathcal S_A=q_{\rm st}(A)\), \(\mathcal S_B=q_{\rm st}(B)\).
Use \(\mathcal S_A\) as the height sites. For an odd state
\(b\in\mathcal S_B\), define its even neighborhood by actual-edge
existence:

\[
 \begin{aligned}
 N_{\mathrm q}(b)=\{a\in\mathcal S_A:\ &
   \text{there are }u\in B,\ v\in A\text{ with }\\
 & q_{\rm st}(u)=b,\ q_{\rm st}(v)=a,\ u\sim v\}.
 \end{aligned}
\]
This definition remains valid when \(\mathcal S_A\cap\mathcal S_B\ne\varnothing\). The random choice at each even image is shared by all its actual even representatives, and every actual odd role reads the entire neighborhood of its odd image. A coarse flip changes at most one bin field and the flag; a fine flip
changes one count field and possibly the field for \(j\); a residual
flip changes one bit. Thus the encoded distance of each edge is bounded,
and two even neighbors of one odd image are within \(D_0=10\).
There are constantly many coarse possibilities and boundedly many
possibilities per fine chunk, giving degree \(O(n)\).
The sets \(\mathcal S_A,\mathcal S_B\) retain their even and odd roles
even when they contain the same encoded point. All neighbors across a
low/high transition form one state at the paired fine-count fringe.
All nonmatching-primary neighbors also form one state, differing only in
the coarse flag. These two facts limit the exceptional tuples below.

An observation descriptor records the distinct ID/type pairs used by the
neighboring even states, together with the coarse and fine keys defining
their laws. Repeated use of one ID/type pair observes the same tuple once;
it does not create another tuple or another hidden scalar. The number of
distinct observations is bounded as follows.
\begin{itemize}
\item At low center severity \(j\), adjacent types observed with at most \(T\) distinct IDs altogether give \(O(T+j+1)\) tuples; as before generic variants are constantly many (same \(t,F\), coarse choices and severity variants), while exceptional adjacent states changing \(F\) or \(t\) number \(O(j+1)\).
\item At high, high observations have constant generic variants; exceptional sign-changing optional variants only matter at \(j\le J+2\), with \(O(J)\) exceptional states. There is at most one adjacent low tuple and at most one nonmatching-primary tuple.
\end{itemize}
Let \(\mathcal D_b\) be the set of descriptors at a fixed odd state
when the allowed IDs form a polynomial-size set. To specify one, list its
at most \(T\) IDs, choose a subset of that list for each of the constantly
many generic types, and give one ID for each exceptional neighboring
state. Therefore
\[
 \log|\mathcal D_b|=O(T\log n+(J+1)\log T).
\]
For unions over histories, use abstract descriptors obtained by renaming
these IDs in \([T]\) and translating the central sign. Their remaining
fine data specify the \(O(j+1)\) coordinates at low states where a flip
can change \(F\), the sign, or the near-fringe status, with constantly
many possible statuses per coordinate. At high states the same description
uses \(O(J)\) coordinates near the interface; away from it the types
have no fine-key data. Consequently their log count per coarse bin and
central sign is
\[
 O(T\log T+(J+1)\log m).
\]
Multiplicities of generic neighboring states do not change the observed
tuple list or its law. The Step~2 type patterns alone have count
\(\exp(O(u_\beta\log m))\) per low severity and constant count at high.

\paragraph{Step 3: deleting a complete tagged tuple.} For an odd state \(b\) fix such an observation configuration for its neighboring even states, and observe the full tuples of the indicated (ID,type) pairs. We form a posterior \(L_b\) of a target single label. We require for good configurations
\[
 L_b\le e^{.16 k}Q_{-c}
\]
for each tuple \(c\) here in the same low/high mode and matching primary, where \(Q_{-c}\) is a reference posterior not using any entries or tag of that tuple. Absolute caps will be \(N\max L_b\le e^{m^{.15}/2}\) at low, and \(N\max L_b\le n^{.05 J}\) at high. Here and below \(n\) is sufficiently large.

\paragraph{Low targets: a hidden scalar.}
Fix the base and all hidden scalars other than
\(\xi=Z_\ell\), where \(\ell=(h_b,t)\). Its prior is \(\pi_\ell\).
The observed data are the complete tagged tuples in the descriptor.
For a candidate \(\xi\), let \(\Theta_b(\xi)\) require the positive
support cases and the following Step~2 tests for each observed type,
evaluated with this candidate: every defining denominator is positive,
\(m_S\ge n^{-\delta_2u_\beta}\), and \(m_S/m_{S-\ell^{\prime}}\ge n^{-\delta_2u_\beta}\)
for every \(\ell^{\prime}\in S\). Include the
Step~1 cap on the target prior. These conditions use the preceding base
and hidden data and the type list, not any tuple tag or entry.
The support and likelihood bounds proved in Steps~1--2 then hold on this
gate over the whole support of each tuple law.

For a tuple of type \(\beta\), use as reference its tag law
\(T_{\beta,-\ell}\), followed by its label sampler with the required
hit to the named variable \(Z_\ell\) omitted. The omitted normalizer
and integrand do not read \(Z_\ell\), so this reference is independent
of the candidate. Other required hits remain, even if their labels happen
to equal the realized value of \(Z_\ell\). Undefined reference cases
use fixed candidate-independent fallbacks.
The tuple likelihood ratio is bounded by
\[
 n^{d_2u_\beta}
 \begin{cases}
 (1+4\epsilon/c_1)^k,&\text{matching primary},\\
 (2/c_1)^k,&\text{nonmatching primary}.
 \end{cases}
\]
The first factor is the tag comparison from Step~2. In the matching case,
adding the omitted hit removes at most \(\epsilon\) mass from a set of
mass at least \(c_1/2\); this gives the first label factor. The second
case uses only the common-hit mass \(c_1/2\).

Write \(g_e(o_e\mid\xi)\) for these tagged-tuple likelihood ratios
and \(Q^{\rm data}\) for the product of their reference laws. Define
\[
 M(o)=\int\Theta_b(\xi)\prod_e g_e(o_e\mid\xi)\,d\pi_\ell(\xi),
 \qquad
 M_{-c}(o)=\int\Theta_b(\xi)\prod_{e\ne c}g_e(o_e\mid\xi)\,d\pi_\ell(\xi).
\]
Normalize the respective integrands to obtain \(L_b\) and \(Q_{-c}\).
The latter omits the complete tagged tuple \(c\). Impose the data tests
\[
 M>0,\qquad M\ge e^{-.02k},\qquad
 M/M_{-c}\ge e^{-.02k}
\]
for the needed deletions. These tests are not added to the candidate gate
of either posterior. Because \(M\) is the true-gated subdensity and
\(\int M_{-c}\,dQ^{\rm data}\le1\), each exception has joint mass at
most \(e^{-.02k}\), integrating the target and tuples in the raw model.

For a matching tuple in the same mode, normalization gives
\[
 \log\frac{L_b}{Q_{-c}}
 \le d_2u_\beta\log n+k\log(1+4\epsilon/c_1)+.02k<.16k.
\]
Here \(u_\beta\le J+1\) and \(d_2\ll\kappa\).
For the absolute cap, there are \(O(T+J)\) observed tuples, of which
at most one has nonmatching primary. Adding all tag costs, that possible
label cost, the prior cap and the denominator cost gives
\(O(J^2\log n)=o(m^{.15})\). Thus
\(N\max L_b\le e^{m^{.15}/2}\). Every supported candidate hits all
observed tuple entries.

\paragraph{High targets: the actual primary in a local joint model.}
Now the target is the named primary \(\xi=P_{h_b}\), so replacing it
also changes some preceding base and hidden laws. Its prior \(\pi^0\)
is the original conditional candidate law given \(V_0\) if it is
\(A_w\), and the original initial law if it is \(V_0\).
Only in the first case is \(V_0\) fixed from the outset.

Define the preceding observation list \(H_{\rm loc}\) from the fixed
tuple descriptor as follows. Include every hidden scalar in the union of
the observed types' \(S\)-lists. Include the parent pair for each
observed type, omitting the target itself. For each included
\(Z_{(h',t')}\), include the base tags indexed by \(C(h')\) and all
parents needed to generate those tags and to define its posterior,
again omitting the target variable itself.
Take the union, counting each primitive variable once. A deterministic
bounded grid neighborhood may pad this list. There are \(O(1)\) coarse
bins and \(O(J)\) hidden scalars: all high types have at most one hidden
key, and only one neighboring low tuple can contribute its larger list.

Given \(\xi\), generate this list by the original local base and
hidden kernels. Unused exterior variables are marginalized. Let
\(g_*(H_{\rm loc}\mid\xi)\) be its density relative to uniform laws
for observed candidate labels and hidden scalars, and \(\Lambda\) for
base tags. Let \(\Theta(\xi,H_{\rm loc})\) require positive support,
all Step~1 cap and deletion tests for the included hidden keys, and
the full Step~2 conjunction just specified for each observed type. This is a finite
conjunction determined before tuple outcomes are observed.
Replacing \(\xi\) recomputes every affected base-tag law, hidden prior,
type law and gate from these inputs. The hidden data are fixed when the
posterior is evaluated, but are integrated when its raw exception
probability is estimated.

For a tuple use reference tag law \(\Lambda\) and the label sampler
omitting the hit to the named variable \(\xi\), with other required
variables held at their observed values. Every tuple imposes that hit;
undefined omission samplers use candidate-independent fallbacks.
Hits to different variable names remain even if their realized values
coincide with \(\xi\). Write \(g_e(o_e\mid\xi,H_{\rm loc})\) for
the conditional tagged-tuple likelihood relative to this reference.
The absolute tag bound in Step~2 gives the same tuple comparisons as in
the low calculation. In particular the possible low tuple has a matching
primary for this hit.

The unnormalized target measure is
\[
 \Theta(\xi,H_{\rm loc})\,g_*(H_{\rm loc}\mid\xi)
 \prod_e g_e(o_e\mid\xi,H_{\rm loc})\,d\pi^0(\xi).
\]
Let \(M\) be its mass and \(L_b\) its normalization. Omitting the
factor \(e=c\) gives \(M_{-c}\) and \(Q_{-c}\).
The hidden likelihood and gate read no part of tuple \(c\), so this
deletion removes its tag and all its entries. After integrating the kept
tuple factors against their conditional references, the deleted density
has mass at most one: the remaining integral is a gated sublaw of
\(H_{\rm loc}\). Hence the tests \(M>0\), \(M\ge e^{-.02k}\), and
\(M/M_{-c}\ge e^{-.02k}\) have raw true-gated exception mass at most
\(e^{-.02k}\) each. This probability integrates \(H_{\rm loc}\) as
well as the target and tuples. As at low targets, the denominator tests
are data tests, not additional candidate gates.

For a matching high tuple \(u_\beta=1\), so the previous deletion
budget gives \(L_b\le e^{.16k}Q_{-c}\). The absolute-width budget has
four contributions:
\[
 \begin{aligned}
 &O(1+d_0\log n) &&\text{from coarse observations},\\
 &O(Jd_1\log n) &&\text{from hidden scalars},\\
 &O((T+J)d_2\log n) &&\text{from tuple tags},\\
 &o(k)+k\log(2/c_1)+.02k &&\text{from labels and the denominator}.
 \end{aligned}
\]
There is only one possible low tuple, so its \(u_\beta=O(J)\) cost is
included in the third line. Matching label costs total
\(O((T+J)k\epsilon/c_1)=o(k)\), and at most one tuple needs the weaker
label factor. The coefficient contributed by the last line to
\(J\log n\) is at most
\(\kappa(\log(2/c_1)+.02)+o(1)<.002\).
Choosing \(d_2\) and the smaller constants with the stated slack makes
the total less than \(.05J\log n\), proving
\(N\max L_b\le n^{.05J}\). Its support again gives common adjacency.

\subsection{Conditioning, center choices, and load estimates}
\label{sec:conditioning-center-choices-and-load-estimates}

We now sample histories on which the preceding local tests hold
simultaneously. At each stage the new avoidance law is normalized at the
fixed entering history; its success probability does not reweight any
earlier draw. All posterior formulas remain those of the raw experiments.
For a raw failure event \(E\) and an exposed history \(H\), the alarm
variable is \(a_E(H)=\Pr_{\rm raw}(E\mid H)\). Its expectation is the
raw failure probability. Thus Markov's inequality changes a bound \(p\)
into \(\Pr(a_E>\sqrt p)\le\sqrt p\). The events \(E\) below include
their stated local true gates.
After the union over deletions, each Step~3 configuration fails with raw
joint mass at most \(e^{-.018k}\) for large \(n\).
After its local factor union, Step~2 failure costs at most
\(n^{-\delta_2u_\beta/2}\).
We seek all Step~1 and Step~2 tests, and, at the final hidden history,
a Step~3 failure probability over independent hypothetical center tuples
at most \(e^{-c_2k}\) for each abstract descriptor.
The deterministic posterior inequalities then follow from their gates.

\emph{Global parent.} Restrict \(V_0\) to values for which every
conditional raw failure probability loses at most half its exponent.
Conditional on \(V_0\), the bin variables are iid and the hidden laws
are invariant under central-sign translation. Only finitely many coarse
shapes and the abstract descriptors above need be tested.
Markov's inequality and the unions over severities give costs bounded by
\[
 \begin{aligned}
 &O(1)n^{-\delta_2/4}
       +\sum_{j\le J}m^{O(j+1)}n^{-\delta_2(j+1)/4},\\
 &\exp\bigl(O(T\log T+(J+1)\log m)-.009k\bigr)
 \end{aligned}
\]
for Step~2 and Step~3 respectively. Both tend to zero for sufficiently
small \(\alpha\); Step~1 has only constantly many shapes.
The tests that hypothetically redraw \(V_0\) still use its original
prior inside their formulas.

\emph{Coarse base.} At a retained \(V_0\), start with the product of
candidate-and-two-tag variables over bins. At each bin forbid failure of
Step~1, or an alarm that the future Step~2 or Step~3 probability exceeds
its next threshold. Conditional future probabilities have the same sign
symmetry and use only a bounded bin neighborhood. The grouped bad-event
probabilities are \(o(1)\) and their dependency degree is bounded, so
Lemma~\ref{lem:conditional-avoidance} supplies the desired law.
The future bounds after this stage are
\(n^{-\delta_2u_\beta/8}\) and \(e^{-.0045k}\).

\emph{Hidden scalars.} At a retained base, start with the product law
of the \(Z\)'s. Forbid Step~2 failure and the alarm that a future
Step~3 failure exceeds the next threshold, halving exponents by Markov
once more. Group by bin and central sign, including all allowed
severities in one group. Every group uses signs within bounded distance
and nearby bins, so the dependency degree is \(m^{O(1)}\).
Even after the pattern unions these probabilities are at most
\(n^{-c_3}\) for some \(c_3>0\) fixed before \(\alpha\).
High tests away from the interface have no hidden-key input and already
satisfy the coarse-stage bound. We may take \(c_2=.001\), independently
of the later choice of \(\alpha\).

For clarity, the hidden-stage pattern union has probability at most
\(n^{-c_3}\), and each group meets at most \(m^C=n^{C\alpha+o(1)}\)
other groups, for an absolute \(C\). Choose \(\alpha\) so that
\(C\alpha<c_3/4\). Charges \(x=n^{-c_3/2}\) then satisfy
\[
 n^{-c_3}\le x(1-x)^{m^C},\qquad m^C x=o(1).
\]
The coarse stage has bounded degree and the same inequality with
charges tending to zero. These checks give both avoidance and the later
cost of deleting constraints touching a fixed local set.
On the admitted entering histories, the gated exceptions are exactly the
failures that need to be excluded. In particular the high-target formula
may redraw \(V_0\) under its original prior, while its alarm after
\(V_0\) has been sampled is simply a function of the true local data.

\paragraph{Local height selection.}

The tuple sampler already imposes every required hit on valid inputs.
There is no optional pool search in this construction, and the even mean
will be controlled by the tag-mixture estimate below. Thus no singleton
removals from Lemma~\ref{lem:constant-broad-high-neighborhood} are needed.
Use its center and height mechanism with linear radius
\(r=\lfloor\rho n\rfloor\), \(\lambda=n^{10}\), and small fixed
\(\rho>0\). Choose the other height parameters to give crowd at most
\(T/2\) and short-truncation error \(o(e^{-2m^{1/5}})\).
Given admitted hidden history, generate the independent tuple arrays.
At each odd state and each pair of consecutive levels, take a maximal
family of disjoint ID sets whose Step~3 descriptors fail. Mark their
union forbidden at the adjacent even sites.

On the polynomial ball-count bounds let \(D_n\) bound the number of
descriptors at that state and level pair. The preceding count gives
\(\log D_n=O(T\log n+(J+1)\log T)<c_2k/2\), after taking
\(\alpha\) small. Disjoint ID sets use independent arrays, and hence
\[
 \Pr(\text{at least }n\text{ disjoint failed sets})
 \le D_n^n e^{-c_2kn}\le e^{-c_2kn/2}.
\]
This beats the \(\exp(O(n))\) union over states and levels.
After forcing one center present with a fixed tuple, at most one member
of a disjoint family contains it; the other \(n-1\) failures give the
same negligible bound. Ball counts change by at most one.

There are therefore fewer than \(n\) marked sets per star and level
pair. Their total loss at an even site-level is \(O(n^2T)\), leaving
at least \(\lambda/3\) eligible centers. Activate centers and make the
uniform eligible-active choice at the long height. The height lemma
implies global success with probability \(1-o(1)\); each odd state sees
at most \(T\) IDs across at most two levels. A failed actual descriptor
would be disjoint from the marked union, contradicting maximality.

One even state makes one shared choice. Each actual odd role uses the
entire neighborhood of its odd state, not only the states represented by
its own actual neighbors. This convention also defines the calculations
under hypothetical replacements. A mark reads only the incident-star
Step~3 tests, their base and hidden inputs, the listed center tuples and
position counts. Eligibility does not consult future alarms, selection
adjustments or global success. The long or short path rule then reads
these eligibility sets and local activations and ties. Invalid choices
produce zero rows, with fixed local fillers where a continued sampling
experiment requires them.

\paragraph{Selected odd posteriors.}
From here on rows are indexed by actual odd roles \(b\); write
\(L_b\) for the posterior at their state \(q_{\rm st}(b)\).
At a high role use \(p_b=L_b\) for the selected tuple descriptor. At a low role, selection itself carries
information about the hidden scalar, so we use the short-height
correction from Lemma~\ref{lem:constant-broad-high-neighborhood}.
Fix the base, all hidden scalars except \(Z_\ell\), and the prospective
center positions. In the raw experiment let \(Z_\ell\sim\pi_\ell\),
regenerate every consulted tuple whose law depends on it, and compute the
short-height selections with their activation and tie randomness.

For a canonical descriptor \(\mathfrak d\), let \(o\) be its complete
tagged-tuple data and \(Q_{\mathfrak d}^{\rm data}\) their product
reference from Step~3. If \(F_\xi(o)\) is the gated tuple likelihood,
then for the raw experiment \(\mathbb P_\xi\) with target fixed at
\(\xi\),
\[
 \mathbb P_\xi(\text{valid proxy presentation of }\mathfrak d,
                  O_{\mathfrak d}\in do)
 =F_\xi(o)a_\xi(o)\,Q_{\mathfrak d}^{\rm data}(do),
 \qquad 0\le a_\xi\le1.
\]
Here \(a_\xi\) integrates the remaining local selection event after
these data are specified. Candidate-dependent unobserved tuples are
regenerated under their changed laws; they are not fixed at their actual
values. Mappings with the same distinct ID/type observations are combined
into one descriptor, so their presentation events are disjoint.
Validity requires legitimate choices, the descriptor size bound, the
true candidate gate, the Step~3 data tests, and position counts at most
\(2\lambda\) in the adjacent balls at all levels. These are the validity
conditions of the previous lemma with its singleton and optional-pool
conditions removed.

Write \(M=\int F_\xi\,d\pi_\ell\) and
\(M^a=\int F_\xi a_\xi\,d\pi_\ell\).
If \(M^a\ge e^{-.02k}M\), use the posterior proportional to
\(F_\xi a_\xi\,d\pi_\ell\); otherwise use \(L_b\).
The adjustment table depends on the fixed inputs and presented data,
not on a separate read of the actual \(Z_\ell\). Use this same table
also for descriptors presented by the long selection. Invalid outcomes
give zero rows. The possible multiplier \(e^{.02k}\) leaves the bounds
\[
 p_b\le e^{.2k}Q_{-c}\quad\text{for matching primary and mode},
 \qquad
 Np_b\le
 \begin{cases}e^{m^{.15}},&b\text{ low},\\n^{.05J},&b\text{ high}.
 \end{cases}
\]
All valid rows remain supported on common neighbors of the tuple entries.

In the raw proxy experiment, the expected mass assigned to a label \(y\)
is at most \(2\pi_\ell(y)\). To see the two contributions separately,
the adjusted normalizer cancels against the actual presentation density,
and disjointness gives
\[
 \mathbb E[p_b^{\rm pr}(y)1_{\rm adjusted}]\le\pi_\ell(y).
\]
On fallback data, \(M^a<e^{-.02k}M\), while each gated observation
likelihood has integral at most one. Hence
\[
 \mathbb E[p_b^{\rm pr}(y)1_{\rm fallback}]
 \le e^{-.02k}(\#\text{descriptors})\pi_\ell(y)=o(\pi_\ell(y)).
\]
The descriptor bound here uses the polynomial position-count gate and
\(O(T\log n+(J+1)\log T)<.01k\). This expectation uses the original
prior \(\pi_\ell\), before avoidance of hidden-scalar alarms.

Let \(\bar p_b^{\rm pr}\) and \(\bar p_b^{\rm long}\) be the rows
averaged over their raw center experiments at a fixed good hidden
history. The proxy computation consults hidden keys at signs within
\(O(\sqrt m)\) of the role's sign: eligibility is determined by incident
stars, and the short tube contains only that many sign changes.
Couple the two rules by the same positions, activations and ties.
On equality of their adjacent path maxima, choices, validity and kernels
agree. Use the short-truncation estimate in its intersection form,
take a union over the \(O(n)\) adjacent even sites, and multiply by the
common cap. This gives
\[
 N\bar p_b^{\rm long}(y)
 \le N\bar p_b^{\rm pr}(y)
       +O(n)e^{m^{.15}}\,o(e^{-2m^{1/5}})+o(1)
 =N\bar p_b^{\rm pr}(y)+o(1).
\]
The last error also includes the exponentially small eligible-size
exceptions. Both long and short rows, before center averaging, use
center inputs only within spatial radius \(r+o(n)\) at fixed history.

\paragraph{Odd loads.}
Three different averages enter the load estimate. We take all averages
over actual odd roles, using zero outside the indicated low or high
subset.

First expose the base. For a fixed label \(y\), use the summand
\(1_{\{b\text{ low}\}}N\pi_{\ell(b)}(y)\), multiplied on raw inputs
by its Step~1 validity indicator. Its cap is \(n^{d_1}\).
Boundary rows contribute at most \(n^{-.05+d_1}=o(1)\).
At an interior key, before coarse avoidance and conditional on \(V_0\),
posterior cancellation bounds the mean by the candidate prior, hence by
\(O(1)\) after normalization. In an \(n\)-th moment, rows with bins
within a fixed grid distance have fraction
\(O((2n^{-.04})^{d_c})\). Their cap cost is negligible.
For separated bins, remove the coarse constraints touching their local
base variables. Lemma~\ref{lem:conditional-avoidance} costs \(O(1)^l\)
for \(l\le n\) terms and then their raw means multiply.
Lemma~\ref{lem:scattered-moments}, Markov's inequality and a union over
labels give a bounded base average simultaneously for all \(y\), with
probability \(1-o(1)\) in the admitted base law.

Next fix such a base and average the low proxy means over the avoided
hidden-scalar law. Sign neighborhoods of radius \(O(\sqrt m)\) have
fraction \(2^{-m+o(m)}\), giving repeat cost
\(n2^{-m+o(m)}e^{m^{.15}}=o(1)\).
For separated signs, remove the hidden-stage constraints touching each
target scalar. With the other keys fixed, integrate those distinct
targets under their original independent priors. The preceding proxy
calculation bounds each mean by \(2N\pi_{\ell(b)}(y)\), and no other
retained function reads that target. Scattered moments give a bounded
average of the proxy means, and the additive comparison gives the same
for long means. High rows contribute deterministically at most
\(n^{-.13J}n^{.05J}=o(1)\).

Finally fix a typical admitted hidden history and take moments over
centers, activations and ties. Long computations at residual distance
more than \(2r+o(n)\) use disjoint primitive randomness. Their comparison
means are the long means just bounded; the absolute row caps handle close
repeats. Thus the normalized average of \(Np_b(y)\) is bounded for
every \(y\), with probability \(1-o(1)\) on center success. In actual
column sums this is
\[
 \sum_{b\in B}p_b(y)
 =\frac{|B|}{N}\operatorname*{avg}_{b\in B}Np_b(y)=o(1),
\]
so all odd columns are at most \(1/10\) for large \(n\).

\paragraph{The additional even history mean.}
For the even reconstruction we need a separate bound on the tag mixtures:
with probability \(1-o(1)\) in the admitted hidden history,
\[
 \operatorname*{avg}_{x\in A}N\sum_iT_{\beta(x)}(i)\mu_i(a)=O(1)
 \qquad\text{for every label }a.
\]
The contribution outside interior low rows with \(j=0\) is bounded
deterministically. Use \(T_\beta\le n^{d_2u_\beta}\Lambda\), tag
balance and the rarity bounds. For example the low rows with \(j\ge1\)
contribute at most
\(O(\sum_{j\ge1}n^{-.13j+d_2(j+1)})\); boundary and high rows are
handled by their corresponding rarity bounds.

For an interior \(j=0\) row define
\[
 f_x(a;H_{\rm base},Z)
 =1_{\rm local\ tests}\,N\sum_iT_{\beta(x)}(i)\mu_i(a),\qquad
 \phi_x(a;H_{\rm base})=\mathbb E_{Z,\rm raw}[f_x\mid H_{\rm base}],
\]
where the indicator requires true \(I_h\in G_\beta\), positivity,
and the Step~2 tests, with value zero on undefined inputs.
Then \(f_x\le K n^{d_2}\), and \(\phi_x\) depends only on a bounded
bin neighborhood. Hold the base except \(I_h\) fixed. The true-gated
subdensity of \(Z_S\) is the normalizer defining \(T_\beta\), so its
cancellation gives
\[
 \mathbb E_{I_h,Z_S\mid H_{\rm base}\setminus I_h}
       [f_x(a;H_{\rm base},Z)]
 \le N\sum_iT_0(i)\mu_i(a).
\]
This uses the fact established in Step~2 that the posterior formula does
not read the actual \(I_h\) a second time.
Now \(P_h=A_w\), and \(T_0\le O(\eta_{A_w})\).
Only after this cancellation average \(A_w\) under its prior dominated
by \(\Pi\). Since \(\mathbb E_\Pi\eta_{A_w}=\Lambda\),
\[
 \mathbb E_{H_{\rm base},\rm raw}[\phi_x\mid V_0]
 \le O\!\left(N\sum_i\Lambda(i)\mu_i(a)\right)=O(1).
\]

Apply scattered moments first to the \(\phi_x\)'s under the coarse
stage. The near-bin fraction is the one in the preceding odd
base estimate, and the cap here is \(K n^{d_2}\); at separated bins remove the touching coarse constraints
before using the raw expectation. This bounds their cube average.
Given such a base, apply scattered moments to the \(f_x\)'s under the
hidden stage. Each uses boundedly many keys at signs within bounded
distance of the row's sign. The same near-sign count applies; removing
the constraints touching separated key lists lets their means integrate
to the \(\phi_x\)'s. Thus the cube average of the tag mixtures is
bounded, although an individual mixture need not have a constant
pointwise cap.

\paragraph{Odd injection and even posterior reconstruction.}

On entering histories with odd columns at most \(1/10\), apply
Lemma~\ref{lem:near-product-injection} with \(d=N\).
The atom caps satisfy \(N^{-1+o(1)}\le N^{-.95}\), the column bound
is below \(.4\), and \(n^2\le N^{.025}\). We obtain an odd injection
whose outputs on at most \(n^2\) queried roles are dominated by
\(1+o(1)\) times their product law.

Fix an actual even role \(v\) and a possible center ID \(c\), with the
type required at its state. Hold the admitted history through \(Z\)
and all other relevant primitive center randomness fixed. The remaining
primitive variable is the complete tagged tuple
\(z=(i,x_1,\ldots,x_k)\) at this ID/type, with raw law \(P_c\).
For odd-neighbor data \(\mathbf y\) define
\[
 F_z(\mathbf y)=1_E\prod_{b\sim v}p_b(y_b),\qquad
 m_c(\mathbf y)=\int F_z(\mathbf y)\,dP_c(z).
\]
All choices and kernel evaluations in this formula are recomputed at
\(z\). The gate \(E\) requires presence and legitimate long selection
of \(c\) at \(v\), valid kernel presentations on the entire odd-state
stars of its actual neighbors, and correct eligible sizes throughout
\(v\)'s height consultation domain. Thus it includes the descriptor
size and position-count bounds. At a fixed hidden history these checks
use center positions, tuples, activations and ties within radius
\(r+o(n)\), also after a hypothetical replacement of \(z\).

For a neighbor \(b\) with matching mode and primary, let
\(\mathcal D_{b,c}\) be the possible descriptors containing this
ID/type among the locally permitted IDs. The lists need not be selected
validly; including all of them makes the index set independent of
\(z\). For each list use its posterior \(Q_{-c}\), which deletes the
tag and all entries of tuple \(c\). Let \(q_b\) be the uniform average
of these reference laws. On undefined cases choose fixed laws
independent of \(z\); if the prerequisite position-count gate fails,
\(E\) is empty. For nonmatching neighbors use the uniform law on \(Y\).
The product \(Q=\prod_{b\sim v}q_b\) is independent of \(z\).
On \(E\), for each matching neighbor,
\[
 p_b\le e^{.2k}|\mathcal D_{b,c}|q_b,
 \qquad
 \log|\mathcal D_{b,c}|
   =O(T\log n+(J+1)\log T)<.1k.
\]
The deleted law uses only other tuple observations and the preceding
hidden data. In particular its type-list gate never reads the removed
tag or entries. Only coarse or fine chunk flips can produce the other
neighbors, so their number is \(o(n^{.4})\). Their absolute caps cost
\(o(kn)\) in total. We have proved
\[
 F_z\le e^{.34kn}Q.
\]

Retain the successful-prehistory indicator while applying the
near-product inequality to the actual odd labels. After their joint law
has been replaced by the product law, remove nonlocal center-success
restrictions and retain \(E\). Under this raw experiment the marginal
data subdensity is \(m_c\). Lemma~\ref{lem:gated-posterior} therefore
bounds the joint mass of \(E\) and
\(m_c=0\) or \(m_c<e^{-.02kn}Q\) by
\((1+o(1))e^{-.02kn}\) for each fixed \(v,c\).
There are \(\exp(O(n))\) such choices, and
\(\exp(O(n)-.02kn)=o(1)\). Thus the probability of a successful
entering history together with any failed requirement for an actually
selected tuple is \(o(1)\).

The posterior proportional to \(F_z\) now has joint density on the \(k\) labels (marginalizing the tag) relative to uniform product at most \(e^{.4kn}\), since the prior width per coordinate given a support tag costs \(\le n^\gamma+\log(2/c_1)=o(n)\). Its average-coordinate marginal is supported on common neighbors: on every positive-mass input the same tuple reference is selected at \(v\) and each valid neighboring output hits \textbf{all} its label entries by the posterior supports at that star. Let \(\mathcal H\) be the labels at which its average marginal exceeds
\(e^{.55n}/N\), so \(|\mathcal H|/N\le e^{-.55n}\).
The joint cap gives
\[
 \Pr(\#\{j:x_j\in\mathcal H\}>.8k)
 \le 2^k e^{.4kn}(|\mathcal H|/N)^{.8k}=o(1),
\]
because \(.8\cdot.55>.4\).
The heavy marginal mass is the expected fraction of tuple entries in
\(\mathcal H\), hence at most \(.8+o(1)\). Restricting the average marginal to
\(\mathcal H^c\) and renormalizing therefore costs a constant factor. Use these row weights at \(v\) with the selected reference only, zero on local failures. Under auxiliary product odd sampling (with any fillers on invalid rows not used on \(E\)), their expected normalized mass at \(a\), given good history and integrating centers as well, is bounded by
\[
 O\Big(N\sum_i T_{\beta(v)}(i)\mu_i(a)\Big).
\]
Before truncation the cancellation, as a measure in \(z\), is
\[
 \int m_c(\mathbf y)\,
       \frac{F_z(\mathbf y)\,dP_c(z)}{m_c(\mathbf y)}\,d\mathbf y
 =dP_c(z)\int F_z(\mathbf y)\,d\mathbf y.
\]
The last integral retains the selection indicator. After averaging the
other raw center inputs it is bounded by the probability of selecting
this fixed tuple on the stated gates. Coordinate averaging therefore
gives the expected empirical mass of \(a\) in the selected tuple. With one forced present center in the ball and its tuple fixed, the probability with local size gates of its valid selection is \(O(1/\lambda)\) on level 0 by uniform eligible active choice, and \(\exp(-n^{\Omega(1)})\) on higher levels by the height bound. Here the tuple is fixed but the other prospective-center positions
remain averaged. The positive-height bound is the position-averaged
supremum over legal pre-activation eligibilities from
Lemma~\ref{lem:height-selection}, not a bound for each fixed position
configuration. Summing presence probabilities thus costs a constant. The unconditional tuple marginal given history costs at most \(2/c_1\) times the indicated tag mixture.

\paragraph{Even loads and completion.}

Let \(w_v(a)\) be the final even row at label \(a\), extended by zero
on local failure. Fix an admitted history through \(Z\) for which the
additional even mean is bounded. For at most \(n\) even roles separated
in residual bits by more than \(2r+o(n)\), keep the successful entering
history while applying the joint odd-output comparison. Their actual
odd-neighbor sets are disjoint, so this comparison replaces all queried
outputs by their product law at cost \(1+o(1)\).
Then remove nonlocal center-success restrictions, retaining each local
gate \(E\). The raw center integrals now factor, and the preceding
cancellation gives comparison means
\(O(N\sum_iT_{\beta(v)}(i)\mu_i(a))\) for the normalized rows
\(Nw_v(a)\). Their cube average is bounded by the additional history
estimate.

For close repeats choose \(\rho\) so that binary entropy in natural
logs satisfies \(H_{\rm bin}(2\rho)<\log2-.55\).
The near fraction then obeys
\[
 f\le\exp(-(\log2-H_{\rm bin}(2\rho))n+o(n)),
 \qquad nf\,O(e^{.55n})=o(1).
\]
Lemma~\ref{lem:scattered-moments}, Markov's inequality and the column
union bound show that the probability of a retained success together
with any normalized even load above a sufficiently large fixed constant
is \(o(1)\). Since \(|A|/N\to0\), a bounded normalized average gives
actual even column sums at most one for large \(n\). The total exceptional joint mass
through the stagewise laws is \(o(1)\), so these events have a common
successful outcome. Hall's theorem supplies distinct even representatives
adjacent to the already distinct odd labels, completing the embedding.

\end{proof}


